\documentclass[10pt,a4paper]{article}

%% paquetes básicos
\usepackage[utf8]{inputenc}
\usepackage[spanish,es-tabla]{babel}
\usepackage{a4wide,color,verbatim,Sweave,url,xargs,amsmath,booktabs,longtable}
\usepackage{graphicx,float}
\usepackage{tikz,xcolor}
\usepackage{pgfplots}
\usepackage{enumitem}

%% bibliotecas TikZ
\usetikzlibrary{automata,positioning,calc,arrows}
\pgfplotsset{compat=1.18}

%% entornos para exams
\newenvironment{question}{\item}{}
\newenvironment{solution}{\comment}{\endcomment}
\newenvironment{answerlist}{\renewcommand{\labelenumii}{(\alph{enumii})}\begin{enumerate}}{\end{enumerate}}

%% comandos para metadatos exams
\newcommand{\exname}[1]{\def\@exname{#1}}
\newcommand{\extype}[1]{\def\@extype{#1}}
\newcommand{\exsolution}[1]{\def\@exsolution{#1}}
\newcommand{\exshuffle}[1]{\def\@exshuffle{#1}}
\newcommand{\exsection}[1]{\def\@exsection{#1}}

%% configuración párrafos
\setlength{\parskip}{0.7ex plus0.1ex minus0.1ex}
\setlength{\parindent}{0em}

\begin{document}
\Sconcordance{concordance:trigonometria_punto_k_v3.tex:trigonometria_punto_k_v3.Rnw:1 %
35 1 1 128 10 1 1 5 1 1 1 5 1 1 1 5 1 1 1 5 15 1 1 2 4 0 1 2 11 1}


\begin{enumerate}


\begin{question}

Un punto K se mueve de un extremo a otro del segmento QT.

El ángulo $\alpha$ y la medida $h$ se relacionan mediante $\sin(\alpha) = \frac{h}{KP}$, de donde se deduce $KP = \frac{h}{\sin(\alpha)} = h \times \csc(\alpha)$.

¿Cuál gráfica muestra las distancias $KP$ cuando K se mueve sobre el segmento?

\begin{answerlist}
\item
\includegraphics[width=5cm]{grafica_A_v3.png}
\item
\includegraphics[width=5cm]{grafica_B_v3.png}
\item
\includegraphics[width=5cm]{grafica_C_v3.png}
\item
\includegraphics[width=5cm]{grafica_D_v3.png}
\end{answerlist}

\end{question}

\begin{solution}

La relación es $KP = h \times \csc(\alpha) = \frac{h}{\sin(\alpha)}$.

Cuando $\alpha$ es pequeño: $\sin(\alpha) \approx 0$, entonces $\csc(\alpha) \to \infty$.
Cuando $\alpha \to 90$ grados: $\sin(\alpha) \to 1$, entonces $\csc(\alpha) \to 1$.

La gráfica debe mostrar una función cosecante decreciente.

La respuesta correcta es B.

\begin{answerlist}
  \item Falso
  \item Verdadero
  \item Falso
  \item Falso
\end{answerlist}
\end{solution}

%% META-INFORMATION
\exname{Trigonometria Punto Movil v3}
\extype{schoice}
\exsolution{0100}
\exshuffle{TRUE}
\exsection{Geometria Metrica}

\end{enumerate}
\end{document}
